% --- Archon physics formalization source begin ---
% archon:physics
% archon:covers PhyXMiniProblems/problem_phyx_mini_0814.lean
% archon:source-report reports/phyx_mini/problem_phyx_mini_0814.source.json

\chapter{Physics problem phyx\_mini\_0814}
\label{ch:PhyXMiniProblems_problem_phyx_mini_0814}

\paragraph{Problem source.}
\#\# Physical scenario

A 2000 kg (19,600 N) elevator with broken cables in a test rig is falling at 4.00 m/s when it contacts a cushioning spring at the bottom of the shaft. The spring is intended to stop the elevator, compressing 2.00 m as it does so as shown in figure. During the motion a safety clamp applies a constant 17,000 N friction force to the elevator.

\#\# Question

What is the necessary force constant \textbackslash{}( k \textbackslash{}) for the spring?

\#\# Answer choices

- A: A: \textbackslash{}( 1.36 \textbackslash{}times 10\textasciicircum{}4 \textbackslash{}text\{ N/m\} \textbackslash{})

- B: B: \textbackslash{}( 1.06 \textbackslash{}times 10\textasciicircum{}4 \textbackslash{}text\{ N/m\} \textbackslash{})

- C: C: \textbackslash{}( 1.06 \textbackslash{}times 10\textasciicircum{}5 \textbackslash{}text\{ N/m\} \textbackslash{})

- D: D: \textbackslash{}( 1.36 \textbackslash{}times 10\textasciicircum{}5 \textbackslash{}text\{ N/m\} \textbackslash{})

\#\# Figure caption (auxiliary; use the image as primary evidence)

The image depicts an elevator system involving two scenarios:

1. **Left Side (Initial Scenario)**:
   - An elevator cabin with mass \textbackslash{}( m = 2000 \textbackslash{}, \textbackslash{}text\{kg\} \textbackslash{}).
   - The cabin is moving downward with velocity \textbackslash{}( v\_1 = 4.00 \textbackslash{}, \textbackslash{}text\{m/s\} \textbackslash{}).
   - An upward force \textbackslash{}( f = 17,000 \textbackslash{}, \textbackslash{}text\{N\} \textbackslash{}) is applied, depicted by a red arrow.
   - Gravitational weight \textbackslash{}( w = mg \textbackslash{}) is shown pulling the cabin downward.
   - A spring is located beneath the cabin.
   - Two points are labeled as Point 1 and Point 2, with a distance of 2.00 meters between them.

2. **Right Side (Final Scenario)**:
   - The elevator cabin is stationary with velocity \textbackslash{}( v\_2 = 0 \textbackslash{}).
   - The cabin rests compressed on the spring.

Both scenarios are encased within a guided shaft, and cables/ropes are visible above the cabins. The transition from the initial to the final scenario involves a change in movement from downward motion to rest due to compression of the spring.

\#\# Dataset metadata

Work and Energy; Multi-Formula Reasoning, Physical Model Grounding Reasoning

\paragraph{Recorded answer/context.}
B: B: \textbackslash{}( 1.06 \textbackslash{}times 10\textasciicircum{}4 \textbackslash{}text\{ N/m\} \textbackslash{})

\paragraph{Figure/image path.}
/root/proposal\_for\_physic/science-mango/phyx\_mini\_run/phyx\_data/test\_image/814.png

\paragraph{Formalization target.}
create a compiling Lean file with sorry bodies at `PhyXMiniProblems/problem\_phyx\_mini\_0814.lean`. The Lean declarations must preserve the physical quantities, dimensions or dimensional roles, figure labels, governing-law hypotheses, and final relation expressed by this problem.
Use Mathlib/Physlib names found through LeanExplore where available. If a physics API is missing, introduce faithful local abstractions rather than scalar placeholder aliases.

\begin{theorem}[Physics formalization target]
\label{thm:physics:phyx_mini_0814:target}
\lean{PhyXMiniProblems.ProblemPhyXMini0814.problem_phyx_mini_0814}
\uses{def:physics:phyx-mini-0814:phyxminiproblems-problemphyxmini0814-springconstantinnewtonspermeter, def:physics:phyx-mini-0814:phyxminiproblems-problemphyxmini0814-elevatorspringsetup, def:physics:phyx-mini-0814:phyxminiproblems-problemphyxmini0814-haselevatorspringproblemdata, def:physics:phyx-mini-0814:phyxminiproblems-problemphyxmini0814-matcheselevatorspringscenario, def:physics:phyx-mini-0814:phyxminiproblems-problemphyxmini0814-matchessuppliedelevatorspringfigure, def:physics:phyx-mini-0814:phyxminiproblems-problemphyxmini0814-hasphysicalelevatorspringparameters, def:physics:phyx-mini-0814:phyxminiproblems-problemphyxmini0814-satisfieselevatorweightlaw, def:physics:phyx-mini-0814:phyxminiproblems-problemphyxmini0814-satisfieselevatorkineticenergylaw, def:physics:phyx-mini-0814:phyxminiproblems-problemphyxmini0814-satisfieslinearspringpotentialenergylaw, def:physics:phyx-mini-0814:phyxminiproblems-problemphyxmini0814-satisfiesconstantforceworklaws, def:physics:phyx-mini-0814:phyxminiproblems-problemphyxmini0814-satisfieselevatorspringworkenergytheorem, def:physics:phyx-mini-0814:phyxminiproblems-problemphyxmini0814-recordeddatasetanswer, def:physics:phyx-mini-0814:phyxminiproblems-problemphyxmini0814-answerchoicereportsspringconstant}
The assigned autoformalize agent should translate this physics problem into Lean declarations in the covered file, with theorem and lemma proof bodies written as `by sorry`.
\end{theorem}
\begin{proof}
This is an autoformalization task, not a proof task. Produce statements that can later be proved by the physics prover without weakening the physical model.
\end{proof}

% --- Archon declaration topology begin ---
\section{Declaration topology}
\begin{definition}[vertical Velocity Dimension]
  \label{def:physics:phyx-mini-0814:phyxminiproblems-problemphyxmini0814-verticalvelocitydimension}
  \lean{PhyXMiniProblems.ProblemPhyXMini0814.verticalVelocityDimension}
  Dimension of signed vertical velocity, L T⁻¹
\end{definition}
\begin{proof}
  This is the declared model definition of vertical Velocity Dimension.
\end{proof}
\begin{definition}[acceleration Dimension]
  \label{def:physics:phyx-mini-0814:phyxminiproblems-problemphyxmini0814-accelerationdimension}
  \lean{PhyXMiniProblems.ProblemPhyXMini0814.accelerationDimension}
  Dimension of acceleration magnitude, L T⁻²
\end{definition}
\begin{proof}
  This is the declared model definition of acceleration Dimension.
\end{proof}
\begin{definition}[force Dimension]
  \label{def:physics:phyx-mini-0814:phyxminiproblems-problemphyxmini0814-forcedimension}
  \lean{PhyXMiniProblems.ProblemPhyXMini0814.forceDimension}
  Dimension of force magnitude, M L T⁻²
\end{definition}
\begin{proof}
  This is the declared model definition of force Dimension.
\end{proof}
\begin{definition}[spring Constant Dimension]
  \label{def:physics:phyx-mini-0814:phyxminiproblems-problemphyxmini0814-springconstantdimension}
  \lean{PhyXMiniProblems.ProblemPhyXMini0814.springConstantDimension}
  Dimension of a linear spring constant, M T⁻² = N/m
\end{definition}
\begin{proof}
  This is the declared model definition of spring Constant Dimension.
\end{proof}
\begin{definition}[Mass Quantity]
  \label{def:physics:phyx-mini-0814:phyxminiproblems-problemphyxmini0814-massquantity}
  \lean{PhyXMiniProblems.ProblemPhyXMini0814.MassQuantity}
  A nonnegative, unit-independent physical mass
\end{definition}
\begin{proof}
  This is the declared model definition of Mass Quantity.
\end{proof}
\begin{definition}[Length Quantity]
  \label{def:physics:phyx-mini-0814:phyxminiproblems-problemphyxmini0814-lengthquantity}
  \lean{PhyXMiniProblems.ProblemPhyXMini0814.LengthQuantity}
  A nonnegative physical distance or spring-compression magnitude
\end{definition}
\begin{proof}
  This is the declared model definition of Length Quantity.
\end{proof}
\begin{definition}[Vertical Velocity Quantity]
  \label{def:physics:phyx-mini-0814:phyxminiproblems-problemphyxmini0814-verticalvelocityquantity}
  \lean{PhyXMiniProblems.ProblemPhyXMini0814.VerticalVelocityQuantity}
  \uses{def:physics:phyx-mini-0814:phyxminiproblems-problemphyxmini0814-verticalvelocitydimension}
  A signed vertical velocity, using the upward-positive convention
\end{definition}
\begin{proof}
  This is the declared model definition of Vertical Velocity Quantity.
\end{proof}
\begin{definition}[Acceleration Magnitude Quantity]
  \label{def:physics:phyx-mini-0814:phyxminiproblems-problemphyxmini0814-accelerationmagnitudequantity}
  \lean{PhyXMiniProblems.ProblemPhyXMini0814.AccelerationMagnitudeQuantity}
  \uses{def:physics:phyx-mini-0814:phyxminiproblems-problemphyxmini0814-accelerationdimension}
  A nonnegative gravitational-acceleration magnitude
\end{definition}
\begin{proof}
  This is the declared model definition of Acceleration Magnitude Quantity.
\end{proof}
\begin{definition}[Force Magnitude Quantity]
  \label{def:physics:phyx-mini-0814:phyxminiproblems-problemphyxmini0814-forcemagnitudequantity}
  \lean{PhyXMiniProblems.ProblemPhyXMini0814.ForceMagnitudeQuantity}
  \uses{def:physics:phyx-mini-0814:phyxminiproblems-problemphyxmini0814-forcedimension}
  A nonnegative physical force magnitude
\end{definition}
\begin{proof}
  This is the declared model definition of Force Magnitude Quantity.
\end{proof}
\begin{definition}[Spring Constant Quantity]
  \label{def:physics:phyx-mini-0814:phyxminiproblems-problemphyxmini0814-springconstantquantity}
  \lean{PhyXMiniProblems.ProblemPhyXMini0814.SpringConstantQuantity}
  \uses{def:physics:phyx-mini-0814:phyxminiproblems-problemphyxmini0814-springconstantdimension}
  A nonnegative linear-spring force constant
\end{definition}
\begin{proof}
  This is the declared model definition of Spring Constant Quantity.
\end{proof}
\begin{definition}[mass In Kilograms]
  \label{def:physics:phyx-mini-0814:phyxminiproblems-problemphyxmini0814-massinkilograms}
  \lean{PhyXMiniProblems.ProblemPhyXMini0814.massInKilograms}
  \uses{def:physics:phyx-mini-0814:phyxminiproblems-problemphyxmini0814-massquantity}
  Coherent-SI kilogram readout of a physical mass
\end{definition}
\begin{proof}
  This is the declared model definition of mass In Kilograms.
\end{proof}
\begin{definition}[length In Meters]
  \label{def:physics:phyx-mini-0814:phyxminiproblems-problemphyxmini0814-lengthinmeters}
  \lean{PhyXMiniProblems.ProblemPhyXMini0814.lengthInMeters}
  \uses{def:physics:phyx-mini-0814:phyxminiproblems-problemphyxmini0814-lengthquantity}
  Coherent-SI metre readout of a distance or compression
\end{definition}
\begin{proof}
  This is the declared model definition of length In Meters.
\end{proof}
\begin{definition}[vertical Velocity In Meters Per Second]
  \label{def:physics:phyx-mini-0814:phyxminiproblems-problemphyxmini0814-verticalvelocityinmeterspersecond}
  \lean{PhyXMiniProblems.ProblemPhyXMini0814.verticalVelocityInMetersPerSecond}
  \uses{def:physics:phyx-mini-0814:phyxminiproblems-problemphyxmini0814-verticalvelocityquantity}
  Coherent-SI signed vertical-velocity readout in metres per second
\end{definition}
\begin{proof}
  This is the declared model definition of vertical Velocity In Meters Per Second.
\end{proof}
\begin{definition}[acceleration In Meters Per Second Squared]
  \label{def:physics:phyx-mini-0814:phyxminiproblems-problemphyxmini0814-accelerationinmeterspersecondsquared}
  \lean{PhyXMiniProblems.ProblemPhyXMini0814.accelerationInMetersPerSecondSquared}
  \uses{def:physics:phyx-mini-0814:phyxminiproblems-problemphyxmini0814-accelerationmagnitudequantity}
  Coherent-SI acceleration-magnitude readout in metres per second squared
\end{definition}
\begin{proof}
  This is the declared model definition of acceleration In Meters Per Second Squared.
\end{proof}
\begin{definition}[force In Newtons]
  \label{def:physics:phyx-mini-0814:phyxminiproblems-problemphyxmini0814-forceinnewtons}
  \lean{PhyXMiniProblems.ProblemPhyXMini0814.forceInNewtons}
  \uses{def:physics:phyx-mini-0814:phyxminiproblems-problemphyxmini0814-forcemagnitudequantity}
  Coherent-SI force-magnitude readout in newtons
\end{definition}
\begin{proof}
  This is the declared model definition of force In Newtons.
\end{proof}
\begin{definition}[spring Constant In Newtons Per Meter]
  \label{def:physics:phyx-mini-0814:phyxminiproblems-problemphyxmini0814-springconstantinnewtonspermeter}
  \lean{PhyXMiniProblems.ProblemPhyXMini0814.springConstantInNewtonsPerMeter}
  \uses{def:physics:phyx-mini-0814:phyxminiproblems-problemphyxmini0814-springconstantquantity}
  Coherent-SI spring-constant readout in newtons per metre
\end{definition}
\begin{proof}
  This is the declared model definition of spring Constant In Newtons Per Meter.
\end{proof}
\begin{definition}[energy In Joules]
  \label{def:physics:phyx-mini-0814:phyxminiproblems-problemphyxmini0814-energyinjoules}
  \lean{PhyXMiniProblems.ProblemPhyXMini0814.energyInJoules}
  Coherent-SI energy or work readout in joules
\end{definition}
\begin{proof}
  This is the declared model definition of energy In Joules.
\end{proof}
\begin{definition}[Elevator Instant]
  \label{def:physics:phyx-mini-0814:phyxminiproblems-problemphyxmini0814-elevatorinstant}
  \lean{PhyXMiniProblems.ProblemPhyXMini0814.ElevatorInstant}
  The two elevator states displayed in the source image
\end{definition}
\begin{proof}
  This is the declared model definition of Elevator Instant.
\end{proof}
\begin{definition}[Figure Point]
  \label{def:physics:phyx-mini-0814:phyxminiproblems-problemphyxmini0814-figurepoint}
  \lean{PhyXMiniProblems.ProblemPhyXMini0814.FigurePoint}
  The two labeled levels whose vertical separation is 2.00 m
\end{definition}
\begin{proof}
  This is the declared model definition of Figure Point.
\end{proof}
\begin{definition}[Vertical Direction]
  \label{def:physics:phyx-mini-0814:phyxminiproblems-problemphyxmini0814-verticaldirection}
  \lean{PhyXMiniProblems.ProblemPhyXMini0814.VerticalDirection}
  Qualitative directions used by the image's arrows and vertical geometry
\end{definition}
\begin{proof}
  This is the declared model definition of Vertical Direction.
\end{proof}
\begin{definition}[Cable Condition]
  \label{def:physics:phyx-mini-0814:phyxminiproblems-problemphyxmini0814-cablecondition}
  \lean{PhyXMiniProblems.ProblemPhyXMini0814.CableCondition}
  Condition of the elevator's supporting cables
\end{definition}
\begin{proof}
  This is the declared model definition of Cable Condition.
\end{proof}
\begin{definition}[Clamp Force Model]
  \label{def:physics:phyx-mini-0814:phyxminiproblems-problemphyxmini0814-clampforcemodel}
  \lean{PhyXMiniProblems.ProblemPhyXMini0814.ClampForceModel}
  Idealization of the safety-clamp force during the two-metre descent
\end{definition}
\begin{proof}
  This is the declared model definition of Clamp Force Model.
\end{proof}
\begin{definition}[Spring Model]
  \label{def:physics:phyx-mini-0814:phyxminiproblems-problemphyxmini0814-springmodel}
  \lean{PhyXMiniProblems.ProblemPhyXMini0814.SpringModel}
  Constitutive model for the cushioning spring
\end{definition}
\begin{proof}
  This is the declared model definition of Spring Model.
\end{proof}
\begin{definition}[Spring Configuration]
  \label{def:physics:phyx-mini-0814:phyxminiproblems-problemphyxmini0814-springconfiguration}
  \lean{PhyXMiniProblems.ProblemPhyXMini0814.SpringConfiguration}
  Qualitative spring states shown before and after compression
\end{definition}
\begin{proof}
  This is the declared model definition of Spring Configuration.
\end{proof}
\begin{definition}[Figure Object]
  \label{def:physics:phyx-mini-0814:phyxminiproblems-problemphyxmini0814-figureobject}
  \lean{PhyXMiniProblems.ProblemPhyXMini0814.FigureObject}
  Physical objects visible in the supplied raster
\end{definition}
\begin{proof}
  This is the declared model definition of Figure Object.
\end{proof}
\begin{definition}[Figure Label]
  \label{def:physics:phyx-mini-0814:phyxminiproblems-problemphyxmini0814-figurelabel}
  \lean{PhyXMiniProblems.ProblemPhyXMini0814.FigureLabel}
  Numerical or symbolic labels printed in the supplied raster
\end{definition}
\begin{proof}
  This is the declared model definition of Figure Label.
\end{proof}
\begin{definition}[Elevator Spring Figure]
  \label{def:physics:phyx-mini-0814:phyxminiproblems-problemphyxmini0814-elevatorspringfigure}
  \lean{PhyXMiniProblems.ProblemPhyXMini0814.ElevatorSpringFigure}
  \uses{def:physics:phyx-mini-0814:phyxminiproblems-problemphyxmini0814-elevatorinstant, def:physics:phyx-mini-0814:phyxminiproblems-problemphyxmini0814-figurepoint, def:physics:phyx-mini-0814:phyxminiproblems-problemphyxmini0814-verticaldirection, def:physics:phyx-mini-0814:phyxminiproblems-problemphyxmini0814-springconfiguration, def:physics:phyx-mini-0814:phyxminiproblems-problemphyxmini0814-figureobject, def:physics:phyx-mini-0814:phyxminiproblems-problemphyxmini0814-figurelabel}
  Literal qualitative and scalar information transcribed from image 814.png. The figure record contains no spring-constant value
\end{definition}
\begin{proof}
  This is the declared model definition of Elevator Spring Figure.
\end{proof}
\begin{definition}[Elevator Spring Setup]
  \label{def:physics:phyx-mini-0814:phyxminiproblems-problemphyxmini0814-elevatorspringsetup}
  \lean{PhyXMiniProblems.ProblemPhyXMini0814.ElevatorSpringSetup}
  \uses{def:physics:phyx-mini-0814:phyxminiproblems-problemphyxmini0814-massquantity, def:physics:phyx-mini-0814:phyxminiproblems-problemphyxmini0814-lengthquantity, def:physics:phyx-mini-0814:phyxminiproblems-problemphyxmini0814-verticalvelocityquantity, def:physics:phyx-mini-0814:phyxminiproblems-problemphyxmini0814-accelerationmagnitudequantity, def:physics:phyx-mini-0814:phyxminiproblems-problemphyxmini0814-forcemagnitudequantity, def:physics:phyx-mini-0814:phyxminiproblems-problemphyxmini0814-springconstantquantity, def:physics:phyx-mini-0814:phyxminiproblems-problemphyxmini0814-elevatorinstant, def:physics:phyx-mini-0814:phyxminiproblems-problemphyxmini0814-verticaldirection, def:physics:phyx-mini-0814:phyxminiproblems-problemphyxmini0814-cablecondition, def:physics:phyx-mini-0814:phyxminiproblems-problemphyxmini0814-clampforcemodel, def:physics:phyx-mini-0814:phyxminiproblems-problemphyxmini0814-springmodel, def:physics:phyx-mini-0814:phyxminiproblems-problemphyxmini0814-elevatorspringfigure}
  Independent physical quantities of the elevator's stopping process. In particular, springConstant is not defined from the recorded answer or from the desired numerical value
\end{definition}
\begin{proof}
  This is the declared model definition of Elevator Spring Setup.
\end{proof}
\begin{definition}[Has Elevator Spring Problem Data]
  \label{def:physics:phyx-mini-0814:phyxminiproblems-problemphyxmini0814-haselevatorspringproblemdata}
  \lean{PhyXMiniProblems.ProblemPhyXMini0814.HasElevatorSpringProblemData}
  \uses{def:physics:phyx-mini-0814:phyxminiproblems-problemphyxmini0814-massinkilograms, def:physics:phyx-mini-0814:phyxminiproblems-problemphyxmini0814-lengthinmeters, def:physics:phyx-mini-0814:phyxminiproblems-problemphyxmini0814-verticalvelocityinmeterspersecond, def:physics:phyx-mini-0814:phyxminiproblems-problemphyxmini0814-forceinnewtons, def:physics:phyx-mini-0814:phyxminiproblems-problemphyxmini0814-elevatorspringsetup}
  Numerical readouts stated in the prose or at the two endpoints. No stiffness readout or answer label occurs here
\end{definition}
\begin{proof}
  This is the declared model definition of Has Elevator Spring Problem Data.
\end{proof}
\begin{definition}[Matches Elevator Spring Scenario]
  \label{def:physics:phyx-mini-0814:phyxminiproblems-problemphyxmini0814-matcheselevatorspringscenario}
  \lean{PhyXMiniProblems.ProblemPhyXMini0814.MatchesElevatorSpringScenario}
  \uses{def:physics:phyx-mini-0814:phyxminiproblems-problemphyxmini0814-elevatorspringsetup}
  Qualitative scenario assumptions stated in the prose
\end{definition}
\begin{proof}
  This is the declared model definition of Matches Elevator Spring Scenario.
\end{proof}
\begin{definition}[Matches Supplied Elevator Spring Figure]
  \label{def:physics:phyx-mini-0814:phyxminiproblems-problemphyxmini0814-matchessuppliedelevatorspringfigure}
  \lean{PhyXMiniProblems.ProblemPhyXMini0814.MatchesSuppliedElevatorSpringFigure}
  \uses{def:physics:phyx-mini-0814:phyxminiproblems-problemphyxmini0814-lengthinmeters, def:physics:phyx-mini-0814:phyxminiproblems-problemphyxmini0814-elevatorspringsetup}
  Primary-image evidence. This records the Point 1/Point 2 geometry, the arrow directions, and the printed labels without imposing a stiffness value
\end{definition}
\begin{proof}
  This is the declared model definition of Matches Supplied Elevator Spring Figure.
\end{proof}
\begin{definition}[Has Physical Elevator Spring Parameters]
  \label{def:physics:phyx-mini-0814:phyxminiproblems-problemphyxmini0814-hasphysicalelevatorspringparameters}
  \lean{PhyXMiniProblems.ProblemPhyXMini0814.HasPhysicalElevatorSpringParameters}
  \uses{def:physics:phyx-mini-0814:phyxminiproblems-problemphyxmini0814-massinkilograms, def:physics:phyx-mini-0814:phyxminiproblems-problemphyxmini0814-lengthinmeters, def:physics:phyx-mini-0814:phyxminiproblems-problemphyxmini0814-accelerationinmeterspersecondsquared, def:physics:phyx-mini-0814:phyxminiproblems-problemphyxmini0814-forceinnewtons, def:physics:phyx-mini-0814:phyxminiproblems-problemphyxmini0814-springconstantinnewtonspermeter, def:physics:phyx-mini-0814:phyxminiproblems-problemphyxmini0814-elevatorspringsetup}
  Positivity conditions for the physical magnitudes in the model
\end{definition}
\begin{proof}
  This is the declared model definition of Has Physical Elevator Spring Parameters.
\end{proof}
\begin{definition}[Satisfies Elevator Weight Law]
  \label{def:physics:phyx-mini-0814:phyxminiproblems-problemphyxmini0814-satisfieselevatorweightlaw}
  \lean{PhyXMiniProblems.ProblemPhyXMini0814.SatisfiesElevatorWeightLaw}
  \uses{def:physics:phyx-mini-0814:phyxminiproblems-problemphyxmini0814-massinkilograms, def:physics:phyx-mini-0814:phyxminiproblems-problemphyxmini0814-accelerationinmeterspersecondsquared, def:physics:phyx-mini-0814:phyxminiproblems-problemphyxmini0814-forceinnewtons, def:physics:phyx-mini-0814:phyxminiproblems-problemphyxmini0814-elevatorspringsetup}
  The weight magnitude obeys w = m g
\end{definition}
\begin{proof}
  This is the declared model definition of Satisfies Elevator Weight Law.
\end{proof}
\begin{definition}[Satisfies Elevator Kinetic Energy Law]
  \label{def:physics:phyx-mini-0814:phyxminiproblems-problemphyxmini0814-satisfieselevatorkineticenergylaw}
  \lean{PhyXMiniProblems.ProblemPhyXMini0814.SatisfiesElevatorKineticEnergyLaw}
  \uses{def:physics:phyx-mini-0814:phyxminiproblems-problemphyxmini0814-massinkilograms, def:physics:phyx-mini-0814:phyxminiproblems-problemphyxmini0814-verticalvelocityinmeterspersecond, def:physics:phyx-mini-0814:phyxminiproblems-problemphyxmini0814-energyinjoules, def:physics:phyx-mini-0814:phyxminiproblems-problemphyxmini0814-elevatorspringsetup}
  Translational kinetic energy obeys K = (1/2) m v² at both endpoints
\end{definition}
\begin{proof}
  This is the declared model definition of Satisfies Elevator Kinetic Energy Law.
\end{proof}
\begin{definition}[Satisfies Linear Spring Potential Energy Law]
  \label{def:physics:phyx-mini-0814:phyxminiproblems-problemphyxmini0814-satisfieslinearspringpotentialenergylaw}
  \lean{PhyXMiniProblems.ProblemPhyXMini0814.SatisfiesLinearSpringPotentialEnergyLaw}
  \uses{def:physics:phyx-mini-0814:phyxminiproblems-problemphyxmini0814-lengthinmeters, def:physics:phyx-mini-0814:phyxminiproblems-problemphyxmini0814-springconstantinnewtonspermeter, def:physics:phyx-mini-0814:phyxminiproblems-problemphyxmini0814-energyinjoules, def:physics:phyx-mini-0814:phyxminiproblems-problemphyxmini0814-elevatorspringsetup}
  Linear-spring potential energy obeys Uₛ = (1/2) k x²
\end{definition}
\begin{proof}
  This is the declared model definition of Satisfies Linear Spring Potential Energy Law.
\end{proof}
\begin{definition}[Satisfies Constant Force Work Laws]
  \label{def:physics:phyx-mini-0814:phyxminiproblems-problemphyxmini0814-satisfiesconstantforceworklaws}
  \lean{PhyXMiniProblems.ProblemPhyXMini0814.SatisfiesConstantForceWorkLaws}
  \uses{def:physics:phyx-mini-0814:phyxminiproblems-problemphyxmini0814-lengthinmeters, def:physics:phyx-mini-0814:phyxminiproblems-problemphyxmini0814-forceinnewtons, def:physics:phyx-mini-0814:phyxminiproblems-problemphyxmini0814-energyinjoules, def:physics:phyx-mini-0814:phyxminiproblems-problemphyxmini0814-elevatorspringsetup}
  For the downward Point 1-to-Point 2 displacement, gravity does positive work and the upward constant clamp force does negative work
\end{definition}
\begin{proof}
  This is the declared model definition of Satisfies Constant Force Work Laws.
\end{proof}
\begin{definition}[Satisfies Elevator Spring Work Energy Theorem]
  \label{def:physics:phyx-mini-0814:phyxminiproblems-problemphyxmini0814-satisfieselevatorspringworkenergytheorem}
  \lean{PhyXMiniProblems.ProblemPhyXMini0814.SatisfiesElevatorSpringWorkEnergyTheorem}
  \uses{def:physics:phyx-mini-0814:phyxminiproblems-problemphyxmini0814-energyinjoules, def:physics:phyx-mini-0814:phyxminiproblems-problemphyxmini0814-elevatorspringsetup}
  Work--energy balance for the elevator-plus-spring system between Point 1 and Point 2. The spring energy is retained explicitly on both sides so the law does not prescribe the unknown stiffness
\end{definition}
\begin{proof}
  This is the declared model definition of Satisfies Elevator Spring Work Energy Theorem.
\end{proof}
\begin{definition}[Spring Constant Answer Choice]
  \label{def:physics:phyx-mini-0814:phyxminiproblems-problemphyxmini0814-springconstantanswerchoice}
  \lean{PhyXMiniProblems.ProblemPhyXMini0814.SpringConstantAnswerChoice}
  Labels attached to the four displayed spring-constant choices
\end{definition}
\begin{proof}
  This is the declared model definition of Spring Constant Answer Choice.
\end{proof}
\begin{definition}[value In Newtons Per Meter]
  \label{def:physics:phyx-mini-0814:phyxminiproblems-problemphyxmini0814-springconstantanswerchoice-valueinnewtonspermeter}
  \lean{PhyXMiniProblems.ProblemPhyXMini0814.SpringConstantAnswerChoice.valueInNewtonsPerMeter}
  \uses{def:physics:phyx-mini-0814:phyxminiproblems-problemphyxmini0814-springconstantanswerchoice}
  Newton-per-metre value printed beside each answer label
\end{definition}
\begin{proof}
  This is the declared model definition of value In Newtons Per Meter.
\end{proof}
\begin{definition}[recorded Dataset Answer]
  \label{def:physics:phyx-mini-0814:phyxminiproblems-problemphyxmini0814-recordeddatasetanswer}
  \lean{PhyXMiniProblems.ProblemPhyXMini0814.recordedDatasetAnswer}
  \uses{def:physics:phyx-mini-0814:phyxminiproblems-problemphyxmini0814-springconstantanswerchoice}
  The answer label recorded in the supplied dataset metadata
\end{definition}
\begin{proof}
  This is the declared model definition of recorded Dataset Answer.
\end{proof}
\begin{definition}[Answer Choice Reports Spring Constant]
  \label{def:physics:phyx-mini-0814:phyxminiproblems-problemphyxmini0814-answerchoicereportsspringconstant}
  \lean{PhyXMiniProblems.ProblemPhyXMini0814.AnswerChoiceReportsSpringConstant}
  \uses{def:physics:phyx-mini-0814:phyxminiproblems-problemphyxmini0814-springconstantinnewtonspermeter, def:physics:phyx-mini-0814:phyxminiproblems-problemphyxmini0814-elevatorspringsetup, def:physics:phyx-mini-0814:phyxminiproblems-problemphyxmini0814-springconstantanswerchoice}
  A displayed choice reports the physical spring constant exactly
\end{definition}
\begin{proof}
  This is the declared model definition of Answer Choice Reports Spring Constant.
\end{proof}
\begin{lemma}[necessary Spring Constant exact]
  \label{lem:physics:phyx-mini-0814:phyxminiproblems-problemphyxmini0814-necessaryspringconstant-exact}
  \lean{PhyXMiniProblems.ProblemPhyXMini0814.necessarySpringConstant_exact}
  \uses{def:physics:phyx-mini-0814:phyxminiproblems-problemphyxmini0814-springconstantinnewtonspermeter, def:physics:phyx-mini-0814:phyxminiproblems-problemphyxmini0814-elevatorspringsetup, def:physics:phyx-mini-0814:phyxminiproblems-problemphyxmini0814-haselevatorspringproblemdata, def:physics:phyx-mini-0814:phyxminiproblems-problemphyxmini0814-satisfieselevatorkineticenergylaw, def:physics:phyx-mini-0814:phyxminiproblems-problemphyxmini0814-satisfieslinearspringpotentialenergylaw, def:physics:phyx-mini-0814:phyxminiproblems-problemphyxmini0814-satisfiesconstantforceworklaws, def:physics:phyx-mini-0814:phyxminiproblems-problemphyxmini0814-satisfieselevatorspringworkenergytheorem}
  The endpoint data and work--energy laws determine (1/2) k (2 m)² = (1/2)(2000 kg)(4 m/s)² + (19600 N)(2 m) - (17000 N)(2 m), so the necessary spring constant is 10600 N/m
\end{lemma}
\begin{proof}
  The conclusion follows from the governing relations and figure readouts named in the statement of necessary Spring Constant exact.
\end{proof}
% --- Archon declaration topology end ---
% --- Archon physics formalization source end ---
